\section{基本概念：Agentic 系统的分类}

\subsection{什么是 Agent？}

「Agent」一词在业界有多种定义：
\begin{itemize}
    \item \textbf{广义}：能独立运行较长时间、使用多种工具完成复杂任务的全自主系统。
    \item \textbf{狭义}：按预定义工作流执行的程序化实现。
\end{itemize}

Anthropic 将所有这些变体统称为 \textbf{Agentic Systems}（代理系统），但在架构层面做出了一个关键区分：

\begin{importantbox}{Workflow vs. Agent 的核心区分}
\begin{itemize}
    \item \textbf{Workflow}（工作流）：LLM 和工具通过\textbf{预定义的代码路径}进行编排。控制流由开发者决定。
    \item \textbf{Agent}（自主代理）：LLM \textbf{动态地自主决定}其处理流程和工具使用。控制权交给模型本身。
\end{itemize}
\end{importantbox}

这一区分的意义在于：Workflow 更可预测、更可控，适合结构化任务；Agent 更灵活，适合开放性问题。两者并非对立，而是复杂度谱系上的不同位置。

\subsection{何时使用（以及何时不使用）Agentic 系统}

Anthropic 给出了一个明确的决策框架：

\begin{knowledgebox}{复杂度递增的决策路径}
\begin{enumerate}
    \item \textbf{首先考虑最简单的方案}：单次 LLM 调用 + Retrieval + In-context Examples 是否已经足够？
    \item \textbf{需要更高可靠性}：使用 Workflow 模式，获得可预测性和一致性。
    \item \textbf{需要灵活性和模型驱动决策}：使用 Agent 模式。
\end{enumerate}
\end{knowledgebox}

\begin{warningbox}{过早引入复杂度的陷阱}
Agentic 系统通常以\textbf{更高的延迟和成本}换取更好的任务表现。在项目初期就引入复杂的 Agent 架构是常见的错误——很多场景下，经过精心优化的单次 LLM 调用完全可以胜任。一定要先评估简单方案的上限，再决定是否引入更复杂的模式。
\end{warningbox}

\subsection{框架的使用建议}

市面上有很多简化 agentic 系统开发的框架，如 Claude Agent SDK、AWS Strands Agents SDK、Rivet（拖拽式 GUI 工作流构建器）等。Anthropic 对框架的建议是：

\begin{itemize}
    \item \textbf{先用原生 API 起步}：很多模式只需几行代码即可实现。
    \item \textbf{理解底层实现}：框架的抽象层可能遮蔽 prompt 和 response 的细节，增加 debug 难度。
    \item \textbf{避免不必要的复杂度}：框架容易诱导开发者在简单方案足够时添加不必要的复杂性。
\end{itemize}

\begin{warningbox}{框架使用的常见错误}
对底层实现的\textbf{错误假设}是客户出错的常见原因。使用框架时，务必理解框架在底层做了什么——它发送了什么 prompt？如何解析 tool call？中间步骤的结果如何传递？不理解这些就无法有效 debug。
\end{warningbox}

\subsection{本章小结}
Agentic 系统分为 Workflow 和 Agent 两大类。Workflow 由代码控制流程，Agent 由模型控制流程。开发时应从简单方案出发，仅在必要时增加复杂度。框架可以帮助快速起步，但要理解底层实现。


